\subsection{Framework}
Basado en la tabla comparativa, se selecciona Flask para el Módulo 4, 
ya que su naturaleza como microframework proporciona flexibilidad y extensibilidad,
permitiendo una integración directa con Python y JSON para gestionar datos de fermentación
enviados por ESP32 vía MQTT, almacenados en MongoDB. Su estructura MVC adaptable optimiza 
el desarrollo de una aplicación pequeña en un solo servidor, facilitando la conexión con React 
mediante APIs ligeras. Aunque requiere configuración manual, esto es aceptable para un sistema compacto. 
En contraste, Django es rígido y orientado a SQL, mientras que FastAPI, enfocado en APIs complejas, excede 
las necesidades actuales. Por ello, Flask es la mejor opción por su simplicidad y escalabilidad.
\caption{Comparación de Frameworks para el módulo 4.}
\label{tab:frameworks}
\begingroup
% Reduce el espacio entre columnas
% Aumenta la altura de las filas
\centering
\small % Usamos fuente pequeña para aprovechar el ancho

\begin{tabularx}{\textwidth}{|l|X|X|X|X|}
\hline
\textbf{Framework} & \textbf{Característica (Modularidad y Extensibilidad)} & \textbf{Estructura} & \textbf{Ventaja} & \textbf{Desventaja} \\
\hline
\textbf{Django} \cite{zotero-item-140}& Posee una \textbf{alta modularidad} gracias a su sistema de "aplicaciones" (\textit{apps}), sin embargo, su extensibilidad es más enfocada a bases de datos SQL. & Sigue un patrón \textbf{MTV} (\textit{Model-Template-View}) rígido que dicta la organización del código; por lo tanto, es un framework "con opinión". & Ofrece una \textbf{solución completa} (\textit{full-stack}) con ORM, \textit{admin panel}, enrutamiento y seguridad integrados, acelerando el desarrollo inicial. & Su \textbf{rigidez estructural} y su diseño intrínseco para SQL complican la integración completa y nativa de MongoDB. \\
\hline
\textbf{Flask} \cite{zotero-item-142}& Al ser un microframework permit una \textbf{excelente modularidad} y \textbf{máxima extensibilidad}, lo que facilita el desarrollador al elegir e integrar librerías externas para cada funcionalidad. & Adopta un patrón \textbf{MVC} (\textit{Model-View-Controller}) flexible o similar, donde el desarrollador define la organización. & Es una framework nativo de Python y JSON, permitiendo una integración directa y de bajo acoplamiento con bases de datos NoSQL basadas en documentos. & Requiere que el desarrollador \textbf{integre y configure} manualmente cada componente o librería externa, incrementando la complejidad de la configuración inicial. \\
\hline
\textbf{FastAPI} \cite{zotero-item-144}& Ofrece \textbf{alta modularidad} al separar la lógica en archivos y rutas (\textit{routers}). Su extensibilidad es simple, gracias al uso de tipado estándar de Python. & Se basa en el concepto de \textbf{funciones de ruta} y \textit{Pydantic models} para la validación de datos. Su estructura está diseñada para la construcción de APIs. & Integración automática y fluida con Python y JSON, gracias al tipado de datos lo que complementa con la generación automática de documentación interactiva(\textit{OpenAPI/ Swagger}). & Está \textbf{estrictamente enfocado en APIs}, por lo tanto, no incluye módulos para \textit{templates} o frontend, lo cual debe ser manejado externamente. \\
\hline
\end{tabularx}
\endgroup
